
\paragraph{Model Collapse.}
Recursive training on synthetic data causes progressive distributional narrowing, a phenomenon termed model collapse.
\citet{shumailov2024collapse} provided the first systematic evidence across multiple model families, showing that tail distributions erode over successive generations.
Theoretical analysis by \citet{dohmatob2024tails} established scaling laws for this tail decay, and subsequent work demonstrated that contamination levels as low as 1\% synthetic data can trigger strong collapse~\cite{dohmatob2025strong}.
The MAD framework~\cite{alemohammad2024mad} formalized the dynamics of self-consuming generative models, confirming that iterative loops converge to degenerate distributions without fresh data injection.
\citet{gerstgrasser2024inevitable} showed that accumulating real data across generations can prevent collapse entirely.
These studies characterize collapse through output distribution metrics, perplexity, or benchmark scores.
Our work complements this line of research: we track collapse through the spectral geometry of learned representations and connect distributional tail decay to measurable contraction in the embedding space.

\paragraph{Spectral Analysis of Representations.}
Spectral properties of neural networks have been studied from several angles.
\citet{martin2021implicit} applied Random Matrix Theory to weight matrices of pretrained networks, showing that heavy-tailed spectral density encodes implicit regularization.
\citet{roy2007effectiverank} defined effective rank as a continuous dimensionality measure based on singular value entropy, which has since become standard in representation analysis~\cite{rankme2023}.
In training dynamics, \citet{he2025spectral} linked spectral collapse of representation covariance matrices to catastrophic forgetting in continual learning.
Closest to our setting, \citet{gu2026sigma} proposed the SIGMA framework for spectral monitoring during recursive synthetic training, with validation on OPT-125M over short generation horizons.
We extend this direction along three axes: systematic tracking over 10 generations, validation at the 410M and 1B scale, and the connection from spectral indicators to closed-loop collapse prevention.

\paragraph{Data Mixing for LLM Training.}
Controlling training data composition is critical for language model quality.
DoReMi~\cite{xie2023doremi} trains a small proxy model to reweight domain mixtures before the main training run.
Data Mixing Laws~\cite{ye2025mixinglaws, data_mixing_laws2024} predict optimal ratios from small-scale experiments so that mixture optimization does not require full-scale retraining.
Both paradigms determine ratios prior to training and hold them fixed throughout.
Our adaptive controller differs fundamentally: it adjusts the real-synthetic mixing ratio online during recursive training, guided by real-time spectral feedback from the model's own representations rather than a priori proxy estimates.